% Time-stamp: <2022-10-30 20:34:56 Romain>
% Romain Lafarguette 2020, https://romainlafarguette.github.io/

%% ---------------------------------------------------------------------------
%% Preamble: Packages and Setup
%% ---------------------------------------------------------------------------
% Class 
\documentclass{beamer}

% Font and encoding
\usepackage[utf8]{inputenc} % Input font
\usepackage[T1]{fontenc} % Output font
\usepackage{lmodern} % Standard LateX font
\usefonttheme{serif} % Standard LateX font

% Maths 
\usepackage{amsfonts, amsmath, mathabx, bm, bbm} % Maths Fonts

% Graphics
\usepackage{graphicx} % Insert graphics
\usepackage{subfig} % Multiple figures in one graphic
\graphicspath{{/../static/img}{/../static/diagrams}}

% Layout
\usepackage{changepage}

% Colors
\usepackage{xcolor}
\definecolor{imfblue}{RGB}{0,76,151} % Official IMF color
\setbeamercolor{title}{fg=imfblue}
\setbeamercolor{frametitle}{fg=imfblue}
\setbeamercolor{structure}{fg=imfblue}
\setbeamercolor{page number in head/foot}{fg=imfblue}
\setbeamerfont{page number in head/foot}{size=\footnotesize}

% Tables
\usepackage{booktabs,rotating,multirow} % Tabular rules and other macros
\usepackage{pdflscape,afterpage} % Landscape mode and afterpage
\usepackage{threeparttable} % Split long tables
\usepackage[font=scriptsize,labelfont=scriptsize,labelfont={color=imfblue}]{caption}

% Import files
\usepackage{import}

% Appendix slides
\usepackage{appendixnumberbeamer} % Manage page numbers for appendix slides

% A few macros: environments
\newenvironment{wideitemize}{\itemize\addtolength{\itemsep}{10pt}}{\enditemize}
\newenvironment{wideenumerate}{\enumerate\addtolength{\itemsep}{10pt}}{\endenumerate}

\newenvironment{extrawideitemize}{\itemize\addtolength{\itemsep}{30pt}}{\enditemize}
\newenvironment{extrawideenumerate}{\enumerate\addtolength{\itemsep}{30pt}}{\endenumerate}


% Define the footer with higher/lower adjustment
\defbeamertemplate{footline}{higher page number}
{
  \hfill
  \usebeamercolor[fg]{page number in head/foot}
  \usebeamerfont{page number in head/foot}  
  \thepage/\inserttotalframenumber\kern1em\vskip2pt %Change xxpt to
                                %lower/higher the footnote
}
\setbeamertemplate{footline}[higher page number]


% Remove navigation symbols and other superfluous elements
\setbeamertemplate{navigation symbols}{}
\setbeameroption{hide notes}
\setbeamertemplate{note page}[plain]
\beamertemplatenavigationsymbolsempty
\hypersetup{pdfpagemode=UseNone} % don't show bookmarks on initial view

% Institute font
\setbeamerfont{institute}{size=\footnotesize}
\DeclareMathSizes{10}{9}{7}{5}  

%% ---------------------------------------------------------------------------
%% Title info
%% ---------------------------------------------------------------------------
\title[Concepts]{Basic Concepts in Time Series}
\author[R. Lafarguette]{Romain Lafarguette}
\institute[IMF]{IMF External Consultant}

\date[STI, 08 Nov 2022]{Singapore Training Institute, 08 November 2022}

\titlegraphic{
    \begin{figure}
    \centering
    \subfloat{{\includegraphics[width=2cm]{../static/img/imf_logo}}}%
    \end{figure}
}

%% ---------------------------------------------------------------------------
%% Title slide
%% ---------------------------------------------------------------------------
\begin{document}

\begingroup
\renewcommand{\insertframenumber}{}
\begin{frame}
  %\addtocounter{framenumber}{-1}
\maketitle
\end{frame}
\endgroup


%% ---------------------------------------------------------------------------
%% Core
%% ---------------------------------------------------------------------------


\begin{frame}
  \frametitle{Random variable}
    \begin{itemize}
    \item A \textbf{random variable} is a measurable function $f: \Omega \mapsto \mathcal{R}$ that assigns to
      a set of outcome $\Omega$ a \textbf{value}, often a real number.
     \item The probability of an outcome is equal to its \textbf{measure} divided by
        the measure of all possible outcomes
     \begin{itemize}
      \item Example: obtaining an even number by rolling a dice: $\{2, 4, 6\}$
      \item Probability to obtain an even number by rolling a dice: $m(\{2, 4,
        6\})/m(\{1, 2, 3, 4, 5, 6\}) = \frac{1}{2}$ \emph{(here, the measure
          simply "counts" the outcomes with equal weights)}
        \end{itemize}     
    \end{itemize}    

    
  \begin{itemize}
  \item Random variables are the "building block" of statistics:
    \begin{itemize}
    \item Random variables are characterized by their distribution (generating
      function, moments, quantiles, etc.)
    \item The behavior of two or more random variables can be characterized by
      their dependence/independence, matrix of variance-covariance, joint
      distribution, etc.
    \item The main theorem of statistics (law of large numbers, central limit
      theorem, etc.) leverages the properties of random variables
    \end{itemize}
  \end{itemize}
\end{frame}



\begin{frame}
  \frametitle{Stochastic Process}

  \begin{itemize}
  \item A stochastic process is a sequence of random variables indexed by time
    ($t$):
    \begin{equation*}
      {\dots, Y_1, Y_2, \dots, Y_t, Y_{t+1}, \dots} = 
    \end{equation*}
    
  \end{itemize}
  

  
\end{frame}


\begin{frame}
  \frametitle{Stationarity}
  
\end{frame}


\begin{frame}
  \frametitle{Ergodicity}
  
\end{frame}


\begin{frame}
  \frametitle{Moments}
  
\end{frame}


\begin{frame}
  \frametitle{Estimator}
  
\end{frame}


\begin{frame}
  \frametitle{Convergence}
  
\end{frame}

\begin{frame}
  \frametitle{Biais}
  
\end{frame}


\begin{frame}
  \frametitle{Efficiency}
  
\end{frame}




%% ---------------------------------------------------------------------------
%% End document
%% ---------------------------------------------------------------------------
\end{document}


%% ---------------------------------------------------------------------------
%% Sample code
%% ---------------------------------------------------------------------------

% \begin{frame}
% \frametitle{Example with a Figure}
%     \makebox[\linewidth]{\includegraphics[width=\paperwidth]{../output/step_007_top_ae_dotplot.pdf}}
% \end{frame}


% \begin{frame}
% \begin{adjustwidth}{-2em}{-2em} % Wider frame 
%   \frametitle{Example with a Wide Table}  
%   \setlength\tabcolsep{4pt}  % default value: 6pt
%   %\footnotesize
%   \centering
%   \input{../output/mytable.tex}\\
% \end{adjustwidth}
% \end{frame}








